\documentclass[11pt]{article}
\usepackage[T1]{fontenc}
\usepackage{lmodern}
\usepackage[a4paper,margin=28mm]{geometry}
\usepackage{amsmath,amssymb,amsthm,microtype}
\usepackage[hidelinks]{hyperref}
\setlength{\emergencystretch}{1.5em}
\newtheorem{theorem}{Theorem}
\newtheorem{lemma}{Lemma}
\newcommand{\ee}{\mathrm e}
\newcommand{\E}{\mathbb E}
\newcommand{\PP}{\mathbb P}
\newcommand{\ind}{\mathbf 1}
\title{Normality of squarefree-weighted series}
\author{Stefan Ringer\thanks{AI tools (GPT 6 Astra and Fable 5.1)
assisted with the development and checking of the proofs.}}
\date{27 September 2026}
\hypersetup{pdftitle={Normality of squarefree-weighted series},
pdfauthor={Stefan Ringer},pdfsubject={Normality and discrepancy; Erdos Problem 259}}
\begin{document}
\maketitle

Erd\H{o}s problem 259 \cite{E259} asks whether
\[
                         \sum_{n\ge1}\frac{n\mu(n)^2}{2^n}
\]
is irrational, where $\mu$ is the M\"obius function, so that $\mu(n)^2$
is the indicator of squarefree integers. Chen and Ruzsa \cite{CR}
proved this, including a stronger statement for infinite subseries.
We establish normality of the full series and extend the result to every
integer denominator base.

A real number is normal to base $B$ if each word of $\ell$ digits occurs
with limiting frequency $B^{-\ell}$. For $x_1,\ldots,x_M\in[0,1)$,
write
\[
 D_M(x_1,\ldots,x_M)=
 \sup_{0\le a<b\le1}
 \left|\frac1M\#\{1\le n\le M:x_n\in[a,b)\}-(b-a)\right|.
\]
Here $\{x\}$ denotes the fractional part of $x$.
We use $\ee(t)=\exp(2\pi it)$,
$\|t\|=\operatorname{dist}(t,\mathbb Z)$, and $[t]_+=\max(t,0)$.
All Fourier frequencies are integers, and logarithms are natural
unless their base is indicated.

\begin{theorem}\label{thm}
For every fixed integer $B\ge2$, the number
\[
                    \alpha_B=\sum_{n\ge1}n\mu(n)^2B^{-n}
\]
is normal to base $B$, with
\[
                D_M((\{B^n\alpha_B\})_{n\le M})
                                      \ll_B(\log\log M)^{-1}.
\]
\end{theorem}

The base $B$ and the number $\alpha_B$ vary together; no assertion
about absolute normality of a fixed $\alpha_B$ is made.
The proof approximates squarefreeness by finitely many divisibility
conditions. On each resulting residue class, the orbit phase is affine
in its index. An average over integer frequencies controls the classes
whose slopes are close to integers.

Write $\vartheta(t)=\sum_{p\le t}\log p$ and
$\pi(t)=\#\{p\le t:p\text{ prime}\}$.
We use the standard consequences of the prime number theorem
\begin{equation}\label{prime}
 \vartheta(y)\sim y,\qquad
 \sum_{p>z}p^{-2}\ll(z\log z)^{-1},\qquad
 \pi(2z)-\pi(z)\gg z/\log z .
\end{equation}

\paragraph{A Fourier estimate for the periodic sieve.}
For large $y$ put
\[
 Q=\prod_{p\le y}p^2,\quad s_y(n)=\prod_{p\le y}\ind_{p^2\nmid n},
 \quad A_r=\sum_{j\ge1}s_y(r+j)B^{-j},\quad
 \widehat\nu_y(m)=\frac1Q\sum_{r\bmod Q}\ee(mA_r).
\]
For an integer $L\ge1$ and a consecutive interval
$W\subseteq\{1,\ldots,L\}$ set
\[
 E_W(m)=\sum_{j\in W}(1-\cos(2\pi mB^{-j})),\quad
 Z_W=\sum_{p\le\min(\sqrt{2L},y)}\lceil |W|/p^2\rceil,\quad
 V=\sum_{\sqrt{2L}<p\le y}p^{-2}.
\]

\begin{lemma}\label{crt}
For every nonzero integer $m$,
\[
 |\widehat\nu_y(m)|
 \le e^{-\frac12V[E_W(m)-2Z_W]_+}
                 +LV^2+\frac{2\pi|m|}{B-1}B^{-L}.
\]
\end{lemma}
\begin{proof}
Take $r$ uniformly modulo $Q$ and truncate $A_r$ at $L$.
The Chinese remainder theorem makes the residues $r\bmod p^2$
independent and uniform. Condition on those for
$p\le\min(\sqrt{2L},y)$, and let $u_j$ indicate that none of these
prime squares divides $r+j$. At most $Z_W$ entries in $W$ vanish.
For each remaining prime, at most one of the $L$ congruences
$r+j\equiv0\pmod{p^2}$ holds. Each has probability $p^{-2}$;
their complement has probability $1-L/p^2>1/2$.

The probability that any position satisfies two such congruences is
at most $LV^2/2$. In the products over these remaining primes, replacing
$\prod_p(1-\ind_{p^2\mid r+j})$ by
$1-\sum_p\ind_{p^2\mid r+j}$ is exact outside this event.
The two resulting complex exponentials have modulus one, so their
averages differ by at most $LV^2$.
The latter average factors, up to a constant of modulus one, as
a product of
\[
 w_p=1-L/p^2+p^{-2}\sum_{j\le L}\ee(-mu_jB^{-j}).
\]
For $\mathcal E_u=\sum_{j\le L}u_j(1-\cos(2\pi mB^{-j}))$,
retaining the pairs involving the complementary outcome in the
nonnegative pair expansion of $1-|w_p|^2$ gives
$1-|w_p|^2\ge2(1-L/p^2)\mathcal E_u/p^2\ge\mathcal E_u/p^2$.
Thus $|w_p|\le e^{-\mathcal E_u/(2p^2)}$.
Since each summand in the energy lies between zero and two,
$\mathcal E_u\ge[E_W(m)-2Z_W]_+$.
Multiply these bounds, average over the conditioned residues, and add
the truncation and replacement errors.
\end{proof}

\paragraph{Exact averaging in an arbitrary integer base.}
For large $y$ set
\[
 \begin{gathered}
 k=\lfloor y\rfloor,\quad J=B^k,\quad T=\lfloor(\log y)^2\rfloor,
 \quad s=\lceil\log_2T\rceil+2,\\
 L=k+\lceil\log_B(TQ)\rceil+\lceil3\log_B(y+2)\rceil,\qquad
 W=\{s+1,\ldots,k\},\quad w=k-s .
 \end{gathered}
\]
Here $s=o(k)$ and $W$ is nonempty eventually. More precisely,
\[
 \frac{w}{L}\longrightarrow\frac{\log B}{\log B+2}>0,
 \qquad w\asymp_B L\asymp_B y.
\]

\begin{lemma}\label{average}
For every fixed $B\ge2$, uniformly over $1\le|\xi|\le T$,
\[
                   \frac1J\sum_{h=1}^J|\widehat\nu_y(h\xi Q)|
                                           \ll_B(\log y)^{-2}.
\]
\end{lemma}
\begin{proof}
Put $A=|\xi|Q$ and $g=\gcd(A,B^k)$.
For sufficiently large $y$, all primes dividing $B$ are at most $y$,
so $v_p(A)=v_p(|\xi|)+2\le\log_2 T+2\le s$ for each such prime.
Consequently $g\mid B^s$.
As $h$ runs from $1$ to $B^k$, $hA\bmod B^k$ is uniform over
the multiples of $g$, each with multiplicity $g$.
Conditional on a residue $r\bmod B^s$ with $0\le r<B^s$ and $g\mid r$, it is therefore
\[
           r+B^sU,\qquad U\text{ uniform on }\{0,\ldots,B^{k-s}-1\}.
\]
The $w$ higher digits are independent and uniform, even when $B$
is composite and $(A,B)>1$.

For every $s<j\le k$, the conditional mean of $\ee(hA/B^j)$ is zero,
as $U$ traverses complete periods of a nontrivial character.
Thus $\E(E_W(hA)\mid r)=w$.
Changing a digit at position $i\ge s$ changes the energy by at most
\[
                   2\pi(B-1)\sum_{j>i}B^{i-j}\le2\pi .
\]
The Azuma--Hoeffding inequality in its bounded-differences form gives,
for $a=19/20$,
\[
 \PP(E_W(hA)<aw\mid r)
       \le\exp\left(-\frac{(1-a)^2w}{2\pi^2}\right)=e^{-w/(800\pi^2)}.
\]
This inequality follows by exposing the independent digits one at a time:
the martingale increment has conditional range at most $2\pi$;
Hoeffding's exponential estimate and Markov's inequality give the display.
Average over $r$ to obtain the same bound without conditioning.

The elementary numerical bound
$\sigma=\sum_pp^{-2}\le\sum_{p\le53}p^{-2}+1/53<0.47$
implies
$Z_W\le\sigma w+\pi(\sqrt{2L})$.
Since $a-2\sigma>1/100$, outside the exponentially small exceptional set,
\[
 E_W(hA)-2Z_W\ge(a-2\sigma)w-2\pi(\sqrt{2L})\ge \delta_B L .
\]
By \eqref{prime}, using
$(\sqrt{2L},2\sqrt{2L}]\subseteq(\sqrt{2L},y]$ eventually,
$V\asymp_B1/(\sqrt L\log L)$.
Also $hA B^{-L}\le(y+2)^{-3}$.
Lemma~\ref{crt} bounds the terms outside the exceptional set by
$e^{-\delta'_B\sqrt y/\log y}+O_B((\log y)^{-2})$.
On the exceptional set use $|\widehat\nu_y|\le1$.
Its contribution is exponentially small in $w$.
Conjugation handles negative $\xi$.
\end{proof}

\paragraph{From Fourier averages to discrepancy.}
The nonnegative Fej\'er kernel
\[
 F_J(t)=1+2\sum_{h=1}^J(1-h/(J+1))\cos(2\pi ht)
       =\frac1{J+1}\left(\frac{\sin(\pi(J+1)t)}{\sin(\pi t)}\right)^2
\]
is interpreted by continuity at integer $t$.
It is at least $4(J+1)/\pi^2$ on $\|t\|\le1/(2(J+1))$.
Thus Lemma~\ref{average} gives
\begin{equation}\label{res}
 \frac1Q\#\{r:\|\xi QA_r\|\le\eta\}\ll_B(\log y)^{-2}
 \quad\left(0<\eta\le\frac1{2(J+1)},\ 1\le|\xi|\le T\right).
\end{equation}
Indeed, average the kernel over $r$, bound its nonconstant coefficients
by their absolute values, and use the Fourier average above.

\begin{proof}[Proof of Theorem~\ref{thm}]
For integer $X$ let $N$ be uniform on $(X,2X]$, and put
\[
 c_B=\frac1{10(1+\log B)},\quad y=c_B\log X,\quad
 H=\lceil4\log_BX\rceil,\quad \eta=\sqrt{Q/X},\quad Y=\log\log X .
\]
Define
\[
 D(N)=\sum_{j\ge1}(N+j)\mu(N+j)^2B^{-j},\qquad
 D_y(N)=\sum_{j\ge1}(N+j)s_y(N+j)B^{-j}.
\]
Exactly $B^N\alpha_B-D(N)\in\mathbb Z$.
An indicator difference within the first $H$ positions requires
$p^2\mid N+j$ for some $p>y$, $j\le H$.
Since $N+j<3X$, the exceptional fraction is at most
\[
 H\sum_{y<p\le\sqrt{3X}}(p^{-2}+X^{-1})\ll_B Y^{-1}.
\]
Outside this event, $|D(N)-D_y(N)|$ is at most
\[
 \tau_X=B^{-H}\left(\frac{2X+H}{B-1}
                          +\frac{B}{(B-1)^2}\right)\ll_B X^{-3}.
\]

For $r=N\bmod Q$ the periodic approximation is
$D_y(N)=NA_r+C_r$, with $C_r=\sum_{j\ge1}j s_y(r+j)B^{-j}$.
Geometric summation in each progression $N=r+Qv$ gives
\[
 |\E\ee(\xi D_y(N))|
 \ll\frac1Q\#\{r:\|\xi QA_r\|\le\eta\}
                       +\frac{Q}{X\eta}+\frac QX .
\]
All parameters are compatible: \eqref{prime} gives
\[
 Q=X^{2c_B+o(1)},\quad J=X^{c_B\log B+o(1)},\quad
 J\eta=X^{c_B(1+\log B)-1/2+o(1)}
                           =X^{-2/5+o(1)}\longrightarrow0.
\]
Also $Q/(X\eta)=\sqrt{Q/X}\to0$ with a fixed power saving.
Therefore \eqref{res} applies, uniformly for $1\le|\xi|\le T$.
Since $T\asymp_B Y^2$, the Erd\H{o}s--Tur\'an inequality \cite{KN}
gives the discrepancy of $\{D_y(N)\}$:
\[
 D_X^{\rm sieve}\ll_B T^{-1}
       +\sum_{\xi=1}^T\frac{|\E\ee(\xi D_y(N))|}{\xi}
                                    \ll_B\frac{\log Y}{Y^2}.
\]
Enlarging or shrinking each test interval by $\tau_X$ gives discrepancy
$O_B(D_X^{\rm sieve}+\tau_X+Y^{-1})=O_B(Y^{-1})$
for $\{D(N)\}$. Thus the exceptional fraction enters only once,
after the Fourier estimate, rather than once per frequency.

Finally decompose a prefix up to $M$ into dyadic blocks, discarding
its first $O(\sqrt M)$ indices. On the remaining blocks
$\log\log X=\log\log M+O(1)$; endpoint rounding costs
$O(\log M/M)$. Weighted averaging proves the stated ordinary rate.
Applying this to base-$B$ digit intervals proves normality.
\end{proof}

The constants and lower thresholds in these asymptotic estimates
may depend on the fixed base $B$.

\begin{thebibliography}{9}
\bibitem{CR} Y.-G. Chen and I. Z. Ruzsa,
\emph{On the irrationality of certain series},
Period. Math. Hungar. \textbf{38} (1999), 31--37,
\href{https://doi.org/10.1023/A:1004742930674}{doi:10.1023/A:1004742930674}.
\bibitem{E259} T. F. Bloom, \emph{Erd\H{o}s Problem 259},
\url{https://www.erdosproblems.com/259}.
\bibitem{KN} L. Kuipers and H. Niederreiter,
\emph{Uniform Distribution of Sequences}, Wiley, 1974.
\end{thebibliography}
\end{document}
